% Ch18 self-study -- "I do / You do" ladder for Zumdahl 18.1-18.4 and
% 18.7. FOURTEEN ladders (56 questions) plus a ten-question mixed set
% = 66 questions, fourteen figures.
%
% SCOPE, verified against the CED PDF on 2026-08-23:
%   IN  -- 18.1-18.4, 18.7 -> CED 9.8, 9.9, 9.10, 9.11
%   OUT -- 18.5-18.6 (batteries, corrosion) and 18.8 (commercial
%          electrolysis) are on TEXTBOOK-MAP's skip list.
%   TWO EXCLUSION STATEMENTS bind this chapter and are honoured from
%   the start rather than patched later:
%     9.8 "Labeling an electrode as positive or negative will not be
%          assessed on the AP Exam."
%     4.7 "The meaning of the terms 'reducing agent' and 'oxidizing
%          agent' will not be assessed on the AP Exam."  -- which is
%          why there is no agent-ranking ladder here, though Zumdahl
%          builds one.
\documentclass{apchem-worksheet}

\apcunitword{Chapter}
\apcunit{18}
\apcunittitle{Electrochemistry}
\apcdoctitle{Self-Study \textbullet\ Chapter 18, I~do / You~do}
\apcsubtitle{Zumdahl \S18.1--18.4, \S18.7 \textbullet\ PDF pp.\ 882--911}

\begin{document}
\apctitleblock

\begin{apcnote}[How to use these notes --- read this first]
Each skill is a \textbf{ladder}: a fully worked example in the
solid-framed box, then \textbf{four} YOUR TURN questions in the dashed
box --- same skill, new numbers, no help. Work all four before checking
the gray \emph{check:} line; a tracker tick needs all four right first
time. A ten-question \textbf{mixed set} follows Ladder 14.

\textbf{What is in and what is out.} \S18.1--18.4 and \S18.7 are
assessed (CED 9.8--9.11). \S18.5--18.6 (batteries, corrosion) and
\S18.8 (commercial electrolysis) are not on the CED and are not treated
here --- skip those pages of Zumdahl.

\textbf{Two things the CED explicitly will not ask you.} First:
\emph{``Labeling an electrode as positive or negative will not be
assessed''} (9.8). Learn \textbf{anode} and \textbf{cathode}, which are
always assessed; the $+$/$-$ signs flip between cell types and the exam
sensibly avoids them. Second: \emph{``The meaning of the terms
`reducing agent' and `oxidizing agent' will not be assessed''} (4.7).
Say ``\ce{Zn} is oxidised'' rather than ``\ce{Zn} is the reducing
agent'' and you are always safe.

\textbf{This chapter is Chapter 17 in different units.} $\Delta G^{\circ}
= -nFE^{\circ}$ is the bridge, and once you have it, every
``is this favoured?'' question you already know how to answer in free
energy you can answer in volts. Ladder 9 is the hinge of the whole
chapter.
\end{apcnote}

% =====================================================================
\section*{\sffamily Ladder 1 \textbullet\ Oxidation numbers}
\zsec{18.1}

Everything in this chapter starts with knowing which atoms changed.
Oxidation numbers are the bookkeeping that tells you.

\begin{center}
\begin{tikzpicture}[font=\sffamily\footnotesize]
  \node[font=\sffamily\small\bfseries] at (5.6,3.6)
    {apply in order --- earlier rules win when two conflict};
  \foreach \y/\num/\rule in {%
      3.0/{1}/{a free element is \textbf{0} --- \ce{Zn}, \ce{O2},
        \ce{Cl2}},
      2.5/{2}/{a monatomic ion equals its \textbf{charge} ---
        \ce{Na+} is $+1$, \ce{S^2-} is $-2$},
      2.0/{3}/{fluorine is always \textbf{$-1$}},
      1.5/{4}/{hydrogen is \textbf{$+1$} (except $-1$ in metal
        hydrides)},
      1.0/{5}/{oxygen is \textbf{$-2$} (except $-1$ in peroxides)},
      0.5/{6}/{the sum equals the \textbf{overall charge} of the
        species}} {
    \node[anchor=west,font=\sffamily\bfseries] at (0.5,\y) {\num.};
    \node[anchor=west] at (1.05,\y) {\rule};}
  \draw[apcgray] (0.4,3.35) -- (11.0,3.35);
  \draw[apcgray] (0.4,0.15) -- (11.0,0.15);
\end{tikzpicture}
\end{center}

\begin{workedexample}[I do: three assignments]
\textbf{(a) Find the oxidation number of Mn in \ce{MnO4-}.}

Oxygen is $-2$ and there are four of them, so they contribute $-8$. The
whole ion carries $-1$, so by rule 6:
\[ x + (-8) = -1 \quad\Longrightarrow\quad x = \mathbf{+7} \]

\textbf{(b) Find the oxidation number of Cr in \ce{Cr2O7^2-}.}

Seven oxygens give $-14$; the ion is $-2$:
\[ 2x - 14 = -2 \quad\Longrightarrow\quad 2x = 12
   \quad\Longrightarrow\quad x = \mathbf{+6} \]
Note the division by 2 --- the answer is per chromium atom, not for
both together.

\textbf{(c) Find the oxidation number of S in \ce{H2SO4}.}

Two hydrogens give $+2$, four oxygens give $-8$, and the molecule is
neutral:
\[ 2 + x - 8 = 0 \quad\Longrightarrow\quad x = \mathbf{+6} \]

\textbf{Why this matters here.} An element whose oxidation number
\emph{rises} has been \textbf{oxidised}; one whose number \emph{falls}
has been \textbf{reduced}. That is the only definition you need, and it
works even when no oxygen or hydrogen is involved.

\textbf{The mnemonic, if you want one.} \textbf{OIL RIG}: Oxidation Is
Loss of electrons, Reduction Is Gain. Losing electrons makes an atom
more positive, which is exactly why its oxidation number rises.
\end{workedexample}

\begin{yourturn}[YOUR TURN 1 --- four questions]
\begin{enumerate}[leftmargin=*,itemsep=0.5em]
  \item Find the oxidation number of N in \ce{NO3-}.
        \workspace[1.6cm]
  \item Find the oxidation number of S in \ce{SO3^2-}.
        \workspace[1.6cm]
  \item Find the oxidation number of Cl in \ce{ClO4-}.
        \workspace[1.6cm]
  \item An element's oxidation number goes from $+2$ to $+4$. Was it
        oxidised or reduced? \workspace[1.6cm]
\end{enumerate}
\selfcheck{(a) $+5$ \quad (b) $+4$ \quad (c) $+7$ \quad (d) oxidised}
\keyonly{\begin{apcanswer}(a) $x - 6 = -1$, so $x = \mathbf{+5}$.
(b) $x - 6 = -2$, so $x = \mathbf{+4}$. (c) $x - 8 = -1$, so
$x = \mathbf{+7}$. (d) \textbf{Oxidised} --- the number \emph{rose},
which means electrons were lost.\end{apcanswer}}
\end{yourturn}

% =====================================================================
\section*{\sffamily Ladder 2 \textbullet\ Half-reactions and balancing
in acid}
\zsec{18.1}

Split the reaction into the part that loses electrons and the part that
gains them, balance each separately, then recombine.

\begin{center}
\begin{tikzpicture}[font=\sffamily\footnotesize]
  \node[font=\sffamily\small\bfseries] at (5.6,3.5)
    {the five steps, in this order};
  \foreach \y/\num/\stp in {%
      2.9/{1}/{split into two \textbf{half-reactions}},
      2.4/{2}/{balance all atoms \emph{except} O and H},
      1.9/{3}/{balance \textbf{O} by adding \ce{H2O}},
      1.4/{4}/{balance \textbf{H} by adding \ce{H+}},
      0.9/{5}/{balance \textbf{charge} by adding \ce{e-}, then scale
        so the electrons cancel}} {
    \node[anchor=west,font=\sffamily\bfseries] at (0.5,\y) {\num.};
    \node[anchor=west] at (1.05,\y) {\stp};}
  \draw[apcgray] (0.4,3.25) -- (11.0,3.25);
  \draw[apcgray] (0.4,0.55) -- (11.0,0.55);
  \node[anchor=west,font=\sffamily\itshape] at (0.5,0.1)
    {O before H before charge --- doing them out of order is what makes
     this feel hard};
\end{tikzpicture}
\end{center}

\begin{workedexample}[I do: permanganate and iron(II) in acid]
\textbf{Balance $\ce{MnO4- + Fe^2+ -> Mn^2+ + Fe^3+}$ in acidic
solution.}

\textbf{Step 1 --- split.} Manganese goes from $+7$ to $+2$ (reduced);
iron goes from $+2$ to $+3$ (oxidised).
\[ \ce{MnO4- -> Mn^2+}
   \qquad
   \ce{Fe^2+ -> Fe^3+} \]

\textbf{Steps 2--4 on the manganese half.} Mn is already balanced. Four
oxygens on the left need four waters on the right; that puts eight
hydrogens on the right, so add eight \ce{H+} on the left:
\[ \ce{MnO4- + 8H+ -> Mn^2+ + 4H2O} \]

\textbf{Step 5 --- charge.} Left is $-1 + 8 = +7$; right is $+2$. Add
five electrons to the left:
\[ \ce{MnO4- + 8H+ + 5e- -> Mn^2+ + 4H2O} \]

\textbf{The iron half is short.} $\ce{Fe^2+ -> Fe^3+ + e-}$.

\textbf{Scale and add.} The manganese half needs five electrons, the
iron half supplies one --- so multiply the iron half by 5:
\[ \ce{MnO4- + 8H+ + 5Fe^2+ -> Mn^2+ + 4H2O + 5Fe^3+} \]

\textbf{Check both, always.} Atoms: 1 Mn, 4 O, 8 H, 5 Fe on each side.
\checkmark\ Charge: left $= -1 + 8 + 10 = +17$; right $= +2 + 0 + 15 =
+17$. \checkmark\ A balanced redox equation must balance
\emph{charge} as well as atoms, and checking it catches almost every
error this ladder can produce.
\end{workedexample}

\begin{yourturn}[YOUR TURN 2 --- four questions]
\begin{enumerate}[leftmargin=*,itemsep=0.5em]
  \item Balance the half-reaction $\ce{Cr2O7^2- -> Cr^3+}$ in acid.
        \workspace[2.0cm]
  \item How many electrons appear in that half-reaction, and on which
        side? \workspace[1.6cm]
  \item Balance $\ce{Zn -> Zn^2+}$ as a half-reaction.
        \workspace[1.4cm]
  \item Why must charge balance as well as atoms?
        \workspace[1.6cm]
\end{enumerate}
\selfcheck{(a) $\ce{Cr2O7^2- + 14H+ + 6e- -> 2Cr^3+ + 7H2O}$ \quad (b)
6, on the left \quad (c) $\ce{Zn -> Zn^2+ + 2e-}$ \quad (d) electrons
are conserved}
\keyonly{\begin{apcanswer}(a) $\ce{Cr2O7^2- + 14H+ + 6e- -> 2Cr^3+ +
7H2O}$. Check charge: left $= -2 + 14 - 6 = +6$; right $= +6$.
\checkmark\ (b) \textbf{Six}, on the \textbf{left} --- chromium is
reduced (from $+6$ to $+3$, two atoms, three electrons each). (c)
$\ce{Zn -> Zn^2+ + 2e-}$. (d) Because \textbf{electrons are
conserved} --- they are real particles, not bookkeeping. If charge does
not balance, electrons have been created or destroyed
somewhere.\end{apcanswer}}
\end{yourturn}

% =====================================================================
\section*{\sffamily Ladder 3 \textbullet\ Balancing in base}
\zsec{18.1}

One extra move at the end: neutralise the \ce{H+} you added.

\begin{workedexample}[I do: the add-hydroxide trick]
\textbf{Balance $\ce{MnO4- + I- -> MnO2 + I2}$ in basic solution.}

\textbf{First balance it as though it were acidic.} The manganese half:
\[ \ce{MnO4- + 4H+ + 3e- -> MnO2 + 2H2O} \]
The iodide half:
\[ \ce{2I- -> I2 + 2e-} \]
Scale by 2 and 3 so six electrons cancel:
\[ \ce{2MnO4- + 8H+ + 6I- -> 2MnO2 + 4H2O + 3I2} \]

\textbf{Now convert to base.} There are eight \ce{H+}, so add eight
\ce{OH-} to \emph{both} sides:
\[ \ce{2MnO4- + 8H+ + 8OH- + 6I- -> 2MnO2 + 4H2O + 3I2 + 8OH-} \]

\textbf{Combine.} On the left, $\ce{8H+ + 8OH-}$ becomes
$\ce{8H2O}$:
\[ \ce{2MnO4- + 8H2O + 6I- -> 2MnO2 + 4H2O + 3I2 + 8OH-} \]

\textbf{Cancel the water that appears on both sides.} Eight on the
left, four on the right --- subtract four from each:
\[ \ce{2MnO4- + 4H2O + 6I- -> 2MnO2 + 3I2 + 8OH-} \]

\textbf{Check.} Charge: left $= -2 + 0 - 6 = -8$; right $= 0 + 0 - 8 =
-8$. \checkmark\ Atoms: Mn 2, I 6, O $8+4 = 12$ on the left against
$4 + 8 = 12$ on the right, H 8 against 8. \checkmark

\textbf{Why the trick works.} Adding the same thing to both sides never
changes an equation. You are simply using $\ce{H+ + OH- -> H2O}$ to
convert an acidic form into an equivalent basic one --- which is
legitimate because in basic solution \ce{H+} is not available in any
meaningful quantity.
\end{workedexample}

\begin{yourturn}[YOUR TURN 3 --- four questions]
\begin{enumerate}[leftmargin=*,itemsep=0.5em]
  \item After acid-balancing, an equation has 6 \ce{H+} on the left.
        How many \ce{OH-} do you add, and where?
        \workspace[1.6cm]
  \item What do the \ce{H+} and \ce{OH-} combine to form?
        \workspace[1.4cm]
  \item Why is it valid to add \ce{OH-} to both sides?
        \workspace[1.6cm]
  \item In basic solution, which species should \emph{not} appear in
        your final answer? \workspace[1.4cm]
\end{enumerate}
\selfcheck{(a) 6, to both sides \quad (b) \ce{H2O} \quad (c) adding
equally to both sides preserves the equation \quad (d) \ce{H+}}
\keyonly{\begin{apcanswer}(a) \textbf{Six}, added to \textbf{both}
sides. (b) \textbf{Water}, \ce{H2O}. (c) Because adding the
\textbf{same quantity to both sides} of an equation leaves it balanced
and equivalent --- it is the same move as adding the same number to both
sides of an algebraic equation. (d) \textbf{\ce{H+}}. A basic solution
has negligible \ce{H+}, so a final answer containing it has not been
converted.\end{apcanswer}}
\end{yourturn}

% =====================================================================
\section*{\sffamily Ladder 4 \textbullet\ The galvanic cell}
\zsec{18.1}

A galvanic (voltaic) cell separates the two half-reactions so the
electrons must travel through a wire --- which is what makes it useful.

\begin{center}
\begin{tikzpicture}[font=\sffamily\footnotesize]
  % left beaker -- anode (fill before stroke)
  \fill[apcgray!22] (0.7,0.5) rectangle (3.3,2.4);
  \draw[very thick] (0.7,3.0) -- (0.7,0.5) -- (3.3,0.5) -- (3.3,3.0);
  \fill[apcgray!70] (1.8,1.1) rectangle (2.1,3.4);
  % label to the SIDE, so the wire can run down to the electrode top
  \node[anchor=east,font=\sffamily\bfseries] at (1.65,3.05) {\ce{Zn}};
  \node at (1.4,0.85) {\ce{Zn^2+}};
  \node[align=center,font=\sffamily\bfseries] at (2.0,-0.15)
    {ANODE\\oxidation};

  % right beaker -- cathode
  \fill[apcgray!22] (6.7,0.5) rectangle (9.3,2.4);
  \draw[very thick] (6.7,3.0) -- (6.7,0.5) -- (9.3,0.5) -- (9.3,3.0);
  \fill[apcgray!70] (7.9,1.1) rectangle (8.2,3.4);
  \node[anchor=west,font=\sffamily\bfseries] at (8.35,3.05) {\ce{Cu}};
  \node at (8.8,0.85) {\ce{Cu^2+}};
  \node[align=center,font=\sffamily\bfseries] at (8.0,-0.15)
    {CATHODE\\reduction};

  % wire
  % wire runs down to y=3.4, the electrode tops -- stopping at 3.9 left
  % a visible gap and the circuit read as disconnected
  \draw[very thick] (1.95,3.4) -- (1.95,4.3) -- (8.05,4.3)
    -- (8.05,3.4);
  \draw[-{Stealth[length=2.5mm]},very thick] (4.4,4.3) -- (5.6,4.3);
  \node[anchor=south,font=\sffamily\bfseries] at (5.0,4.4)
    {\ce{e-} flow};

  % salt bridge
  \fill[apcgray!45] (3.3,2.0) rectangle (6.7,2.4);
  \draw[very thick] (3.3,2.0) rectangle (6.7,2.4);
  \node[font=\sffamily\bfseries] at (5.0,2.2) {salt bridge};
  \node[align=center,font=\sffamily\footnotesize] at (5.0,1.5)
    {carries ions to keep\\each beaker neutral};

  \node[font=\sffamily\footnotesize] at (5.0,-1.0)
    {electrons always flow \textbf{anode $\to$ cathode} through the
     wire --- never through the salt bridge};
\end{tikzpicture}
\end{center}

\begin{apcnote}[AP scope: learn anode and cathode, not $+$ and $-$]
CED topic 9.8's exclusion: \emph{``Labeling an electrode as positive or
negative will not be assessed on the AP Exam.''} Good --- the signs
reverse between galvanic and electrolytic cells and cause more
confusion than they resolve. What \textbf{is} assessed, and never
changes: \textbf{oxidation always happens at the anode, reduction
always at the cathode}, in every cell of every kind. Learn that pair and
you never need the signs.
\end{apcnote}

\begin{workedexample}[I do: reading a cell]
\textbf{Take the zinc--copper cell above. Identify each part and say
why it is there.}

\textbf{The anode is zinc}, because zinc is oxidised:
$\ce{Zn -> Zn^2+ + 2e-}$. The electrode slowly dissolves and the solution
gains \ce{Zn^2+}.

\textbf{The cathode is copper}, because copper(II) is reduced:
$\ce{Cu^2+ + 2e- -> Cu}$. Copper metal plates onto the electrode, which
gains mass, and the blue colour of the solution fades as \ce{Cu^2+} is
consumed.

\textbf{Electrons flow from anode to cathode} through the external
wire. They are produced where oxidation happens and consumed where
reduction happens, so the direction is forced.

\textbf{The salt bridge exists because of charge build-up.} Without it,
the anode compartment would accumulate positive charge (\ce{Zn^2+}
entering solution) and the cathode compartment negative charge
(\ce{Cu^2+} leaving). Within seconds that imbalance would stop the
reaction. The bridge lets spectator ions migrate --- anions toward the
anode, cations toward the cathode --- keeping both sides neutral.

\textbf{Two memory hooks worth having.} \textbf{An~Ox} and
\textbf{Red~Cat}: \emph{An}ode--\emph{Ox}idation, \emph{Red}uction--%
\emph{Cat}hode. And both words in each pair start with the same kind of
letter --- \textbf{a}node and \textbf{o}xidation are both vowels,
\textbf{c}athode and \textbf{r}eduction both consonants.
\end{workedexample}

\begin{yourturn}[YOUR TURN 4 --- four questions]
\begin{enumerate}[leftmargin=*,itemsep=0.5em]
  \item At which electrode does oxidation occur, in \emph{any} cell?
        \workspace[1.4cm]
  \item Which way do electrons flow through the external wire?
        \workspace[1.4cm]
  \item What would happen without a salt bridge?
        \workspace[1.8cm]
  \item Does the anode of a zinc--copper cell gain or lose mass?
        Explain. \workspace[1.6cm]
\end{enumerate}
\selfcheck{(a) the anode \quad (b) anode to cathode \quad (c) charge
builds up and the reaction stops \quad (d) loses --- \ce{Zn} dissolves}
\keyonly{\begin{apcanswer}(a) The \textbf{anode} --- in every cell,
galvanic or electrolytic. (b) From \textbf{anode to cathode}: they are
produced by oxidation and consumed by reduction. (c) \textbf{Charge
would build up} --- positive in the anode compartment, negative in the
cathode compartment --- and within moments that imbalance would stop
electron flow and the reaction. (d) It \textbf{loses} mass: zinc metal
is oxidised to \ce{Zn^2+} and dissolves into the
solution.\end{apcanswer}}
\end{yourturn}

% =====================================================================
\section*{\sffamily Ladder 5 \textbullet\ Cell notation}
\zsec{18.1}

A one-line shorthand for a whole cell diagram. It has a fixed grammar.

\begin{center}
\begin{tikzpicture}[font=\sffamily\footnotesize]
  \node[font=\sffamily\small] at (5.6,3.2)
    {$\ce{Zn(s)}\;\big|\;\ce{Zn^2+}(aq)\;\big\|\;
     \ce{Cu^2+}(aq)\;\big|\;\ce{Cu(s)}$};
  \node[font=\sffamily\bfseries] at (2.0,2.55) {ANODE side};
  \node[font=\sffamily\bfseries] at (9.0,2.55) {CATHODE side};
  \node[font=\sffamily\bfseries] at (5.6,2.55) {salt bridge};
  \draw[apcgray] (0.5,2.2) -- (11.0,2.2);
  \foreach \y/\sym/\mean in {%
      1.7/{$|$}/{a \textbf{phase boundary} --- solid meets solution},
      1.15/{$\|$}/{the \textbf{salt bridge} --- separates the two
        half-cells},
      0.6/{order}/{\textbf{anode on the left}, cathode on the right,
        always}} {
    \node[anchor=west,font=\sffamily\bfseries] at (0.7,\y) {\sym};
    \node[anchor=west] at (2.2,\y) {\mean};}
  \draw[apcgray] (0.5,0.25) -- (11.0,0.25);
  \node[anchor=west] at (0.7,-0.2)
    {read it left to right and you are reading the electrons' journey};
\end{tikzpicture}
\end{center}

\begin{workedexample}[I do: writing and reading notation]
\textbf{(a) Write the notation for a cell in which
$\ce{Mg -> Mg^2+ + 2e-}$ and $\ce{Ag+ + e- -> Ag}$.}

Magnesium is oxidised, so it is the anode and goes on the left; silver
is reduced, so it is the cathode and goes on the right:
\[ \ce{Mg(s)} \;\big|\; \ce{Mg^2+}(aq) \;\big\|\;
   \ce{Ag+}(aq) \;\big|\; \ce{Ag(s)} \]

\textbf{(b) Read $\ce{Fe(s)}|\ce{Fe^2+}\|\ce{Cu^2+}|\ce{Cu(s)}$.}

Iron is on the left, so iron is the anode and is oxidised:
$\ce{Fe -> Fe^2+ + 2e-}$. Copper is on the right, so \ce{Cu^2+} is
reduced: $\ce{Cu^2+ + 2e- -> Cu}$. Overall:
$\ce{Fe + Cu^2+ -> Fe^2+ + Cu}$.

\textbf{When a species is not a solid, you need an inert electrode.}
For a half-cell like $\ce{Fe^3+}/\ce{Fe^2+}$ --- both species dissolved
--- there is no metal to serve as the electrode, so platinum is used and
written into the notation:
\[ \ce{Pt(s)} \;\big|\; \ce{Fe^2+}, \ce{Fe^3+} \;\big\|\; \ldots \]
The comma separates two species in the \emph{same} phase; the single
bar is reserved for a genuine phase boundary.

\textbf{The one rule that prevents most errors:} anode on the left.
Everything else follows from it --- and if you write the cell backwards,
your $E^{\circ}_{\text{cell}}$ in Ladder 7 will come out negative for a
reaction that is actually favoured.
\end{workedexample}

\begin{yourturn}[YOUR TURN 5 --- four questions]
\begin{enumerate}[leftmargin=*,itemsep=0.5em]
  \item In cell notation, which electrode is written on the left?
        \workspace[1.4cm]
  \item What does the double bar represent? \workspace[1.4cm]
  \item Write the notation for a cell with $\ce{Ni -> Ni^2+ + 2e-}$ and
        $\ce{Ag+ + e- -> Ag}$. \workspace[1.8cm]
  \item Why is an inert platinum electrode sometimes needed?
        \workspace[1.8cm]
\end{enumerate}
\selfcheck{(a) the anode \quad (b) the salt bridge \quad (c)
$\ce{Ni}|\ce{Ni^2+}\|\ce{Ag+}|\ce{Ag}$ \quad (d) when no solid is
present to conduct}
\keyonly{\begin{apcanswer}(a) The \textbf{anode}. (b) The \textbf{salt
bridge} separating the two half-cells. (c)
$\ce{Ni(s)}\,|\,\ce{Ni^2+}(aq)\,\|\,\ce{Ag+}(aq)\,|\,\ce{Ag(s)}$.
(d) Because some half-cells contain \textbf{no solid conductor} --- for
example $\ce{Fe^2+}/\ce{Fe^3+}$, where both species are dissolved.
Platinum provides a surface for electron transfer without taking part in
the reaction.\end{apcanswer}}
\end{yourturn}

% =====================================================================
\section*{\sffamily Ladder 6 \textbullet\ Standard reduction potentials}
\zsec{18.2}

Every half-reaction is tabulated as a \textbf{reduction}, measured
against the hydrogen electrode, which is defined as exactly zero.

\begin{center}
\begin{tikzpicture}[font=\sffamily\footnotesize]
  \draw[-{Stealth[length=3mm]},very thick] (1.6,0.4) -- (1.6,4.2);
  \node[anchor=south,font=\sffamily\bfseries] at (1.6,4.25)
    {more easily reduced};
  \foreach \y/\hr/\ev in {%
      3.7/{$\ce{F2 + 2e- -> 2F-}$}/{$+2.87$},
      3.2/{$\ce{Cl2 + 2e- -> 2Cl-}$}/{$+1.36$},
      2.7/{$\ce{Ag+ + e- -> Ag}$}/{$+0.80$},
      2.2/{$\ce{Cu^2+ + 2e- -> Cu}$}/{$+0.34$},
      1.7/{$\ce{2H+ + 2e- -> H2}$}/{$0.00$},
      1.2/{$\ce{Fe^2+ + 2e- -> Fe}$}/{$-0.44$},
      0.7/{$\ce{Zn^2+ + 2e- -> Zn}$}/{$-0.76$}} {
    \fill (1.6,\y) circle (0.07);
    \node[anchor=west] at (1.9,\y) {\hr};
    \node[anchor=west] at (6.6,\y) {\ev\ V};}
  \node[anchor=west,font=\sffamily\bfseries] at (8.1,1.7)
    {$\leftarrow$ defined as 0};
  \node[anchor=west,align=left,font=\sffamily\footnotesize]
    at (8.1,3.2)
    {a \textbf{more positive} $E^{\circ}$\\
     means that species is\\
     reduced more readily};
\end{tikzpicture}
\end{center}

\begin{apctrap}
\textbf{When you reverse a half-reaction, flip the sign of
$E^{\circ}$ --- but do \emph{not} multiply it when you scale the
coefficients.} Doubling $\ce{Zn -> Zn^2+ + 2e-}$ does not double its
potential. $E^{\circ}$ is an \emph{intensive} property, like density or
temperature: it does not depend on how much you have. This is the one
place where electrochemistry breaks the Hess's-law habit of scaling
everything.
\end{apctrap}

\begin{workedexample}[I do: what the table is telling you]
\textbf{Everything in the table is written as a reduction, so a
positive value means ``this happens readily as written''.}

\textbf{Fluorine, $+2.87$ V.} The most positive value on any standard
table --- \ce{F2} is reduced with enormous willingness, which is why
elemental fluorine is so ferociously reactive.

\textbf{Zinc, $-0.76$ V.} Negative, so $\ce{Zn^2+ -> Zn}$ is
\emph{unfavourable} as written. Run it backwards ---
$\ce{Zn -> Zn^2+ + 2e-}$, oxidation --- and the sign flips to $+0.76$ V,
which is favourable. Zinc prefers to be oxidised.

\textbf{Hydrogen, exactly $0.00$ V.} Not measured but \emph{defined}:
every other potential is measured relative to the standard hydrogen
electrode. Same logical move as $\Delta H_{f}^{\circ}$ of an element
being zero in Chapter 6 --- a reference point, not a physical fact.

\textbf{How to use the table for a spontaneity question.} The species
higher on the table (more positive) is reduced; the one lower down runs
in reverse and is oxidised. In a zinc--copper cell, copper is higher, so
\ce{Cu^2+} is reduced and zinc is oxidised. That single reading gives
you the anode, the cathode, and the direction --- before you compute
anything.

\textbf{The terminology caution.} You may see this table described as
ranking ``oxidising agents''. The CED's topic 4.7 exclusion states that
\emph{``the meaning of the terms `reducing agent' and `oxidising agent'
will not be assessed''}, so phrase your answers as ``\ce{Cu^2+} is
reduced'' rather than ``\ce{Cu^2+} is the oxidising agent''. The
chemistry is identical; only the wording is safer.
\end{workedexample}

\begin{yourturn}[YOUR TURN 6 --- four questions]
Use the table above.
\begin{enumerate}[leftmargin=*,itemsep=0.5em]
  \item Which is more readily reduced, \ce{Ag+} or \ce{Cu^2+}?
        \workspace[1.4cm]
  \item Give $E^{\circ}$ for $\ce{Zn -> Zn^2+ + 2e-}$.
        \workspace[1.4cm]
  \item Give $E^{\circ}$ for $\ce{2Zn -> 2Zn^2+ + 4e-}$.
        \workspace[1.6cm]
  \item Why is the hydrogen half-reaction exactly $0.00$ V?
        \workspace[1.6cm]
\end{enumerate}
\selfcheck{(a) \ce{Ag+} \quad (b) $+0.76$ V \quad (c) still $+0.76$ V
\quad (d) it is a defined reference}
\keyonly{\begin{apcanswer}(a) \textbf{\ce{Ag+}} --- $+0.80$ V against
copper's $+0.34$ V. (b) $\mathbf{+0.76}$ \textbf{V}: reversing the
tabulated reduction flips the sign. (c) \textbf{Still $+0.76$ V.}
Scaling the coefficients does \emph{not} scale $E^{\circ}$ --- it is an
intensive property. (d) Because it is \textbf{defined} as the
reference point; every other potential is measured relative to the
standard hydrogen electrode. It is a convention, like
$\Delta H_{f}^{\circ} = 0$ for an element.\end{apcanswer}}
\end{yourturn}

% =====================================================================
\section*{\sffamily Ladder 7 \textbullet\ Calculating
$E^{\circ}_{\text{cell}}$}
\zsec{18.2}

\[ E^{\circ}_{\text{cell}}
   = E^{\circ}_{\text{cathode}} - E^{\circ}_{\text{anode}} \]

where both values are read straight off the table as \emph{reductions}
--- no sign flipping needed.

\begin{workedexample}[I do: the Daniell cell, two ways]
\textbf{Find $E^{\circ}_{\text{cell}}$ for the zinc--copper cell.}

\textbf{Method 1 --- cathode minus anode.} Copper is reduced (cathode,
$+0.34$ V); zinc is oxidised (anode, tabulated as $-0.76$ V):
\[ E^{\circ}_{\text{cell}} = 0.34 - (-0.76)
   = \mathbf{+1.10\;\text{V}} \]
Note you subtract the \emph{tabulated} anode value, minus sign and all.
Do not flip it first --- the formula already does that.

\textbf{Method 2 --- flip and add, if you prefer.} Write the oxidation
explicitly and flip its sign:
\[ \ce{Zn -> Zn^2+ + 2e-} \qquad +0.76\;\text{V} \]
\[ \ce{Cu^2+ + 2e- -> Cu} \qquad +0.34\;\text{V} \]
\[ E^{\circ}_{\text{cell}} = 0.76 + 0.34 = +1.10\;\text{V} \]
Same answer. Use whichever you find harder to get wrong --- but never
mix them, which is where sign errors come from.

\textbf{The electrons must cancel, and they do so without scaling
$E^{\circ}$.} Both half-reactions here involve two electrons, so nothing
needs multiplying. Even if they had not --- say a silver half-cell
needing $\times 2$ --- you would scale the \emph{equation} but leave
$E^{\circ}$ untouched.

\textbf{Sanity check.} A positive $E^{\circ}_{\text{cell}}$ means the
cell as written runs on its own, which is exactly what a battery does.
If your answer is negative, either you have the anode and cathode
swapped, or the reaction genuinely needs to be driven (Ladder 13).
\end{workedexample}

\begin{yourturn}[YOUR TURN 7 --- four questions]
Use $E^{\circ}$: \ce{Ag+}/\ce{Ag} $= +0.80$, \ce{Cu^2+}/\ce{Cu} $=
+0.34$, \ce{Fe^2+}/\ce{Fe} $= -0.44$, \ce{Zn^2+}/\ce{Zn} $= -0.76$ V.
\begin{enumerate}[leftmargin=*,itemsep=0.5em]
  \item Find $E^{\circ}_{\text{cell}}$ for a \ce{Zn}/\ce{Ag} cell.
        \workspace[1.8cm]
  \item Find $E^{\circ}_{\text{cell}}$ for a \ce{Fe}/\ce{Cu} cell.
        \workspace[1.8cm]
  \item In the \ce{Fe}/\ce{Cu} cell, which metal is the anode?
        \workspace[1.4cm]
  \item If you doubled every coefficient, what happens to
        $E^{\circ}_{\text{cell}}$? \workspace[1.6cm]
\end{enumerate}
\selfcheck{(a) $+1.56$ V \quad (b) $+0.78$ V \quad (c) iron \quad (d)
nothing --- it is unchanged}
\keyonly{\begin{apcanswer}(a) Silver is higher, so it is the cathode:
$0.80 - (-0.76) = \mathbf{+1.56}$ \textbf{V}. (b) Copper is higher, so
it is the cathode: $0.34 - (-0.44) = \mathbf{+0.78}$ \textbf{V}. (c)
\textbf{Iron} --- it is lower on the table, so it is oxidised. (d)
\textbf{Nothing.} $E^{\circ}$ is intensive and does not scale with the
amount of substance.\end{apcanswer}}
\end{yourturn}

% =====================================================================
\section*{\sffamily Ladder 8 \textbullet\ Spontaneity from
$E^{\circ}_{\text{cell}}$}
\zsec{18.3}

\begin{center}
\begin{tikzpicture}[font=\sffamily\footnotesize]
  \draw[-{Stealth[length=3mm]},very thick] (0.6,2.2) -- (10.9,2.2);
  \node[font=\sffamily\bfseries] at (5.7,2.7) {$E^{\circ} = 0$};
  \fill (5.7,2.2) circle (0.11);
  \draw[densely dashed,apcgray] (5.7,1.92) -- (5.7,2.15);
  \node[font=\sffamily\bfseries] at (2.6,1.75) {$E^{\circ} > 0$};
  \node[align=center] at (2.6,1.0)
    {thermodynamically\\\textbf{favoured} --- a galvanic cell};
  \node[align=center] at (5.7,1.0) {at\\equilibrium};
  \node[font=\sffamily\bfseries] at (8.8,1.75) {$E^{\circ} < 0$};
  \node[align=center] at (8.8,1.0)
    {must be \textbf{driven} ---\\an electrolytic cell};
  \draw[apcgray] (0.5,0.45) -- (11.0,0.45);
  \node[anchor=west] at (0.6,0.0)
    {note the sign convention is \emph{opposite} to $\Delta G$: positive
     $E^{\circ}$ and negative $\Delta G^{\circ}$ both mean favoured};
\end{tikzpicture}
\end{center}

\begin{apctrap}
\textbf{The signs run opposite ways, and that is the trap.} Favoured
means $\Delta G^{\circ} < 0$ but $E^{\circ} > 0$. The minus sign in
$\Delta G^{\circ} = -nFE^{\circ}$ is what flips them, and mixing the two
conventions is the commonest error in Unit 9. When in doubt, remember a
battery: it works on its own, and it has a \emph{positive} voltage.
\end{apctrap}

\begin{workedexample}[I do: predicting whether a reaction happens]
\textbf{Will copper metal dissolve in a solution of \ce{Zn^2+}?}

\textbf{Write the proposed reaction.} If it happened, copper would be
oxidised and zinc(II) reduced:
\[ \ce{Cu + Zn^2+ -> Cu^2+ + Zn} \]

\textbf{Assign the electrodes as proposed.} Zinc is being reduced, so it
would be the cathode ($-0.76$ V); copper is oxidised, so the anode
($+0.34$ V):
\[ E^{\circ}_{\text{cell}} = -0.76 - 0.34 = -1.10\;\text{V} \]

\textbf{Negative, so it does not happen.} Copper will not dissolve in
zinc(II) solution. And note the answer is exactly the negative of
Ladder 7's $+1.10$ V --- because it is the same reaction written
backwards.

\textbf{Now the reverse, which does happen.} Drop zinc metal into
\ce{Cu^2+} solution and $E^{\circ} = +1.10$ V: the zinc dissolves, copper
plates onto it, and the blue colour fades. This is a standard
demonstration and a standard exam question.

\textbf{The shortcut, once you trust the table.} A metal will displace
any metal ion \emph{below} it in the table. Zinc is below copper, so
zinc displaces \ce{Cu^2+}; copper does not displace \ce{Zn^2+}. That
single reading answers most ``will this reaction occur'' questions
without arithmetic.
\end{workedexample}

\begin{yourturn}[YOUR TURN 8 --- four questions]
Use the $E^{\circ}$ values from Ladder 7.
\begin{enumerate}[leftmargin=*,itemsep=0.5em]
  \item Will zinc metal displace \ce{Ag+} from solution? Show the
        potential. \workspace[1.8cm]
  \item Will silver metal displace \ce{Cu^2+}? Show the potential.
        \workspace[1.8cm]
  \item What is the sign of $\Delta G^{\circ}$ when $E^{\circ}$ is
        positive? \workspace[1.4cm]
  \item Why do a positive $E^{\circ}$ and a negative $\Delta G^{\circ}$
        mean the same thing? \workspace[1.8cm]
\end{enumerate}
\selfcheck{(a) yes, $+1.56$ V \quad (b) no, $-0.46$ V \quad (c)
negative \quad (d) the minus sign in $\Delta G = -nFE$}
\keyonly{\begin{apcanswer}(a) \textbf{Yes}: $0.80 - (-0.76) =
\mathbf{+1.56}$ \textbf{V}, favoured. (b) \textbf{No}: silver would be
oxidised and \ce{Cu^2+} reduced, giving $0.34 - 0.80 = \mathbf{-0.46}$
\textbf{V}. (c) \textbf{Negative}. (d) Because $\Delta G^{\circ} =
-nFE^{\circ}$, and $n$ and $F$ are both positive --- the
\textbf{minus sign} converts a positive potential into a negative free
energy. They are the same statement in different
units.\end{apcanswer}}
\end{yourturn}

% =====================================================================
\section*{\sffamily Ladder 9 \textbullet\ $\Delta G^{\circ} =
-nFE^{\circ}$}
\zsec{18.3}

The hinge of the chapter --- the equation that makes electrochemistry a
branch of Chapter 17.

\[ \Delta G^{\circ} = -n F E^{\circ}
   \qquad
   F = \SI{96485}{\coulomb\per\mole} \]

$n$ is the number of moles of electrons transferred in the balanced
equation.

\begin{center}
\begin{tikzpicture}[font=\sffamily\footnotesize]
  \node[font=\sffamily\small\bfseries] at (5.6,3.0)
    {three descriptions of one fact};
  \node[font=\sffamily\bfseries] at (1.9,2.2) {$E^{\circ}_{\text{cell}}$};
  \node[font=\sffamily\bfseries] at (5.6,2.2) {$\Delta G^{\circ}$};
  \node[font=\sffamily\bfseries] at (9.3,2.2) {$K$};
  \draw[{Stealth[length=2mm]}-{Stealth[length=2mm]},thick]
    (2.6,2.2) -- (4.9,2.2);
  \node[font=\sffamily\footnotesize,anchor=south] at (3.75,2.3)
    {$-nF$};
  \draw[{Stealth[length=2mm]}-{Stealth[length=2mm]},thick]
    (6.3,2.2) -- (8.6,2.2);
  \node[font=\sffamily\footnotesize,anchor=south] at (7.45,2.3)
    {$-RT\ln$};
  \draw[apcgray] (0.5,1.7) -- (11.0,1.7);
  \foreach \y/\e/\g/\kk in {%
      1.2/{$E^{\circ} > 0$}/{$\Delta G^{\circ} < 0$}/{$K > 1$},
      0.7/{$E^{\circ} = 0$}/{$\Delta G^{\circ} = 0$}/{$K = 1$},
      0.2/{$E^{\circ} < 0$}/{$\Delta G^{\circ} > 0$}/{$K < 1$}} {
    \node at (1.9,\y) {\e};
    \node at (5.6,\y) {\g};
    \node at (9.3,\y) {\kk};}
\end{tikzpicture}
\end{center}

\begin{workedexample}[I do: volts to kilojoules and back]
\textbf{(a) The zinc--copper cell has $E^{\circ} = +1.10$ V and
transfers 2 mol of electrons. Find $\Delta G^{\circ}$.}
\[ \Delta G^{\circ} = -nFE^{\circ}
   = -(2)(96485)(1.10) \]
\[ = \SI{-212000}{\joule} = \mathbf{\SI{-212}{\kilo\joule}} \]

Negative, as a positive $E^{\circ}$ requires --- and large, which is why
this cell is a usable battery.

\textbf{Units check.} A volt is a joule per coulomb, so
$\text{mol} \times \text{C/mol} \times \text{J/C}$ leaves joules.
Convert to kilojoules at the end; forgetting to leaves an answer a
thousand times too large.

\textbf{(b) A reaction transferring 2 mol of electrons has
$\Delta G^{\circ} = \SI{-96.5}{\kilo\joule}$. Find $E^{\circ}$.}

Rearrange:
\[ E^{\circ} = \frac{-\Delta G^{\circ}}{nF}
   = \frac{96500}{(2)(96485)}
   = \mathbf{+0.500\;\text{V}} \]
Note the numbers were chosen so $F$ nearly cancels --- exam items often
are.

\textbf{Finding $n$ is the step people skip.} It is the number of
electrons that actually cancel when you combine the half-reactions ---
\emph{not} the number in either half alone. For
$\ce{2Al + 3Cu^2+ -> 2Al^3+ + 3Cu}$, aluminium gives up 3 each and
copper takes 2 each, so the balanced equation transfers
$2 \times 3 = 6$ moles of electrons: $n = 6$.
\end{workedexample}

\begin{yourturn}[YOUR TURN 9 --- four questions]
Use $F = \SI{96485}{\coulomb\per\mole}$.
\begin{enumerate}[leftmargin=*,itemsep=0.5em]
  \item $E^{\circ} = +0.50$ V and $n = 2$. Find $\Delta G^{\circ}$ in
        kJ. \workspace[2.0cm]
  \item $E^{\circ} = -0.25$ V and $n = 1$. Find $\Delta G^{\circ}$ and
        say whether it is favoured. \workspace[2.0cm]
  \item For $\ce{2Al + 3Cu^2+ -> 2Al^3+ + 3Cu}$, what is $n$?
        \workspace[1.6cm]
  \item Why does a positive $E^{\circ}$ always give a negative
        $\Delta G^{\circ}$? \workspace[1.6cm]
\end{enumerate}
\selfcheck{(a) \SI{-96.5}{\kilo\joule} \quad (b)
\SI{+24.1}{\kilo\joule}, not favoured \quad (c) 6 \quad (d) the minus
sign, with $n$ and $F$ positive}
\keyonly{\begin{apcanswer}(a) $-(2)(96485)(0.50) = \SI{-96485}{\joule}
= \mathbf{\SI{-96.5}{\kilo\joule}}$. (b) $-(1)(96485)(-0.25) =
\SI{+24121}{\joule} = \mathbf{\SI{+24.1}{\kilo\joule}}$ ---
\textbf{not favoured}. (c) Aluminium loses 3 electrons each and there
are 2 of them, so $n = \mathbf{6}$; equivalently 3 copper ions gain 2
each. (d) Because $n$ and $F$ are both \textbf{positive} constants, so
the \textbf{minus sign} in $-nFE^{\circ}$ inverts the sign of
$E^{\circ}$.\end{apcanswer}}
\end{yourturn}

% =====================================================================
\section*{\sffamily Ladder 10 \textbullet\ $E^{\circ}$ and $K$}
\zsec{18.3}

Combine Ladder 9 with Chapter 17's $\Delta G^{\circ} = -RT\ln K$ and the
free energy drops out:

\[ E^{\circ} = \frac{0.0592}{n}\,\log K
   \qquad\text{at \SI{298}{\kelvin}} \]

\begin{workedexample}[I do: how far does a battery reaction go?]
\textbf{Find $K$ for the zinc--copper cell, where $E^{\circ} = 1.10$ V
and $n = 2$.}

Rearrange for $\log K$:
\[ \log K = \frac{n E^{\circ}}{0.0592}
          = \frac{(2)(1.10)}{0.0592}
          = \frac{2.20}{0.0592} = 37.2 \]
\[ K = 10^{37.2} \approx \mathbf{10^{37}} \]

\textbf{That is an enormous number}, and it should be. A cell potential
of 1.10 V is large, so the reaction runs essentially to completion ---
which is why the cell delivers current until one reactant is exhausted
rather than fading to an equilibrium mixture.

\textbf{Notice how little voltage it takes.} Just 1.10 V corresponds to
$K \approx 10^{37}$. The logarithm compresses ferociously: every
$0.0592/n$ volts is one power of ten in $K$. A cell with
$E^{\circ} = 0.10$ V and $n = 2$ has $\log K = 3.4$, so $K \approx
10^{3}$ --- still product-favoured, but a billion billion billion times
less so.

\textbf{Where the $0.0592$ comes from.} It is $2.303\,RT/F$ evaluated at
\SI{25}{\celsius} --- that is \SI{298.15}{\kelvin}, which is why the
constant is $0.0592$ and not the $0.0591$ you get from a rounded
\SI{298}{\kelvin}. It is a temperature-specific shortcut, so state the
assumption if a question gives you a different temperature.

\textbf{Sanity check, always.} Positive $E^{\circ}$ gives
$\log K > 0$, so $K > 1$ and products are favoured. Negative
$E^{\circ}$ gives $K < 1$. The three-way consistency of Ladder 9's
figure means any two of $E^{\circ}$, $\Delta G^{\circ}$ and $K$ let you
check the third.
\end{workedexample}

\begin{yourturn}[YOUR TURN 10 --- four questions]
Use $E^{\circ} = (0.0592/n)\log K$ at \SI{298}{\kelvin}.
\begin{enumerate}[leftmargin=*,itemsep=0.5em]
  \item $E^{\circ} = +0.0592$ V and $n = 1$. Find $\log K$ and $K$.
        \workspace[1.8cm]
  \item $E^{\circ} = +0.30$ V and $n = 2$. Find $\log K$.
        \workspace[1.8cm]
  \item If $E^{\circ}$ is negative, is $K$ greater or less than 1?
        \workspace[1.4cm]
  \item Why is the constant $0.0592$ only valid at
        \SI{298}{\kelvin}? \workspace[1.6cm]
\end{enumerate}
\selfcheck{(a) $\log K = 1$, $K = 10$ \quad (b) $\log K = 10.1$ \quad
(c) less than 1 \quad (d) it contains $T$}
\keyonly{\begin{apcanswer}(a) $\log K = (1)(0.0592)/0.0592 =
\mathbf{1}$, so $K = \mathbf{10}$. (b) $\log K = (2)(0.30)/0.0592 =
\mathbf{10.1}$, so $K \approx 10^{10}$. (c) \textbf{Less than 1} ---
a negative $E^{\circ}$ gives a negative $\log K$. (d) Because
$0.0592$ is $2.303\,RT/F$ \textbf{evaluated at \SI{298}{\kelvin}} ---
it contains the temperature, so it changes if the temperature
does.\end{apcanswer}}
\end{yourturn}

% =====================================================================
\section*{\sffamily Ladder 11 \textbullet\ The Nernst equation}
\zsec{18.4}

Real cells rarely sit at standard conditions. The Nernst equation
corrects $E^{\circ}$ for the concentrations you actually have.

\[ E = E^{\circ} - \frac{0.0592}{n}\,\log Q
   \qquad\text{at \SI{298}{\kelvin}} \]

\begin{center}
\begin{tikzpicture}[font=\sffamily\footnotesize]
  \node[font=\sffamily\small\bfseries] at (5.6,2.85)
    {the same $Q$-versus-$K$ logic as Chapter 13, in volts};
  \foreach \y/\cond/\eff in {%
      2.2/{$Q < K$ (excess reactant)}/{$E > E^{\circ}$ --- more
        driving force},
      1.65/{$Q = K$}/{$E = 0$ --- the cell is \textbf{dead}},
      1.1/{$Q > K$ (excess product)}/{$E < 0$ --- it would run
        backwards}} {
    \node[anchor=west] at (0.8,\y) {\cond};
    \node[anchor=west] at (5.6,\y) {\eff};}
  \draw[apcgray] (0.6,2.6) -- (11.0,2.6);
  \draw[apcgray] (0.6,0.75) -- (11.0,0.75);
  \node[anchor=west] at (0.8,0.3)
    {a flat battery is a cell that has reached equilibrium --- $E = 0$
     is the definition of dead};
\end{tikzpicture}
\end{center}

\begin{workedexample}[I do: a cell away from standard conditions]
\textbf{For the zinc--copper cell ($E^{\circ} = 1.10$ V, $n = 2$), find
$E$ when $[\ce{Zn^2+}] = \SI{0.10}{\Molar}$ and $[\ce{Cu^2+}] =
\SI{1.0}{\Molar}$.}

\textbf{Write $Q$ for the cell reaction} $\ce{Zn + Cu^2+ -> Zn^2+ +
Cu}$. Solids are omitted, exactly as in Chapter 13:
\[ Q = \frac{[\ce{Zn^2+}]}{[\ce{Cu^2+}]}
     = \frac{0.10}{1.0} = 0.10 \]

\textbf{Substitute.}
\[ E = 1.10 - \frac{0.0592}{2}\log(0.10)
     = 1.10 - (0.0296)(-1) \]
\[ = 1.10 + 0.0296 = \mathbf{+1.13\;\text{V}} \]

\textbf{Higher than standard, and that makes sense.} There is less
product (\ce{Zn^2+}) than standard and a full complement of reactant, so
the reaction has further to go and pushes harder. Le Ch\^{a}telier, read
as a voltage.

\textbf{The two things to get right.} First, \textbf{solids and pure
liquids do not appear in $Q$} --- so \ce{Zn} and \ce{Cu} metal are
absent. Second, watch the sign: $\log$ of a number below 1 is negative,
and subtracting a negative \emph{raises} $E$. Sign slips here are the
main source of wrong answers.

\textbf{And the limiting case.} As the cell runs, \ce{Zn^2+} builds up
and \ce{Cu^2+} is consumed, so $Q$ rises. When $Q$ finally reaches $K$,
the correction term exactly cancels $E^{\circ}$ and $E = 0$. The battery
is flat --- not because it ran out of chemicals, but because it reached
equilibrium.
\end{workedexample}

\begin{yourturn}[YOUR TURN 11 --- four questions]
Use $E^{\circ} = 1.10$ V, $n = 2$ for the \ce{Zn}/\ce{Cu} cell.
\begin{enumerate}[leftmargin=*,itemsep=0.5em]
  \item Find $E$ when $[\ce{Zn^2+}] = \SI{1.0}{\Molar}$ and
        $[\ce{Cu^2+}] = \SI{0.010}{\Molar}$. \workspace[2.2cm]
  \item Does raising $[\ce{Cu^2+}]$ raise or lower $E$? Why?
        \workspace[1.8cm]
  \item What is $E$ when the cell reaches equilibrium?
        \workspace[1.4cm]
  \item Which species are omitted from $Q$, and why?
        \workspace[1.6cm]
\end{enumerate}
\selfcheck{(a) $+1.04$ V \quad (b) raises it --- more reactant \quad
(c) 0 \quad (d) solids and pure liquids}
\keyonly{\begin{apcanswer}(a) $Q = 1.0/0.010 = 100$, so
$E = 1.10 - (0.0296)\log(100) = 1.10 - (0.0296)(2) = 1.10 - 0.0592 =
\mathbf{+1.04}$ \textbf{V}. (b) \textbf{Raises} it: more reactant means
a smaller $Q$, a more negative $\log Q$, and therefore a larger $E$.
(c) \textbf{Zero} --- a cell at equilibrium is a flat battery.
(d) \textbf{Solids and pure liquids}, because their concentrations are
fixed and cannot change --- the same rule as every equilibrium
expression since Chapter 13.\end{apcanswer}}
\end{yourturn}

% =====================================================================
\section*{\sffamily Ladder 12 \textbullet\ Concentration cells}
\zsec{18.4}

The strangest cell in the course: identical electrodes, identical
chemistry on both sides, and it still produces a voltage.

\begin{center}
\begin{tikzpicture}[font=\sffamily\footnotesize]
  % dilute side
  \fill[apcgray!22] (0.7,0.5) rectangle (3.3,2.3);
  \draw[very thick] (0.7,2.9) -- (0.7,0.5) -- (3.3,0.5) -- (3.3,2.9);
  \fill[apcgray!70] (1.8,1.0) rectangle (2.1,3.3);
  \foreach \p in {(1.2,0.9),(2.6,1.2)} {\fill \p circle (0.07);}
  \node[align=center] at (2.0,-0.2)
    {\textbf{dilute} \ce{Cu^2+}\\\SI{0.010}{\Molar}};
  \node[font=\sffamily\bfseries] at (2.0,3.55) {\ce{Cu}};

  % concentrated side
  \fill[apcgray!22] (6.7,0.5) rectangle (9.3,2.3);
  \draw[very thick] (6.7,2.9) -- (6.7,0.5) -- (9.3,0.5) -- (9.3,2.9);
  \fill[apcgray!70] (7.9,1.0) rectangle (8.2,3.3);
  \foreach \p in {(7.1,0.85),(7.5,1.4),(8.6,0.95),(8.9,1.7),
                  (7.3,1.95),(8.5,2.0)} {\fill \p circle (0.07);}
  \node[align=center] at (8.0,-0.2)
    {\textbf{concentrated} \ce{Cu^2+}\\\SI{1.0}{\Molar}};
  \node[font=\sffamily\bfseries] at (8.0,3.55) {\ce{Cu}};

  \fill[apcgray!45] (3.3,1.9) rectangle (6.7,2.3);
  \draw[very thick] (3.3,1.9) rectangle (6.7,2.3);
  \node[font=\sffamily\footnotesize] at (5.0,2.1) {salt bridge};
  % Only the short symbol fits in the 3.4cm gap between the beakers;
  % the explanation ran over both beaker walls, so it moved to the
  % caption block below the figure.
  \node[font=\sffamily\footnotesize] at (5.0,1.3) {$E^{\circ} = 0$};
  \node[font=\sffamily\footnotesize] at (5.0,-1.0)
    {both half-cells run the same reaction, so $E^{\circ} = 0$};
  \node[font=\sffamily\footnotesize] at (5.0,-1.5)
    {the \textbf{dilute} side is the anode --- the cell runs to even
     the two concentrations out};
\end{tikzpicture}
\end{center}

\begin{workedexample}[I do: voltage from concentration alone]
\textbf{A copper concentration cell has \SI{0.010}{\Molar} \ce{Cu^2+}
on one side and \SI{1.0}{\Molar} on the other. Find $E$.}

\textbf{Start with $E^{\circ} = 0$.} Both half-reactions are
$\ce{Cu^2+ + 2e- -> Cu}$, so the cathode and anode values are identical
and cancel exactly. Whatever voltage appears comes entirely from the
Nernst correction.

\textbf{Decide which side is the anode.} The cell will act to
\emph{reduce} the difference, so the dilute side must gain \ce{Cu^2+}
--- meaning copper dissolves there. Dilute side is the anode; the
concentrated side is the cathode, where \ce{Cu^2+} plates out.

\textbf{Write $Q$ as (anode concentration) over (cathode
concentration)}, since the anode species is the product:
\[ Q = \frac{0.010}{1.0} = 0.010 \]

\textbf{Apply Nernst with $n = 2$.}
\[ E = 0 - \frac{0.0592}{2}\log(0.010)
     = -(0.0296)(-2) = \mathbf{+0.0592\;\text{V}} \]

\textbf{Small but real.} Concentration cells produce tens of millivolts,
not volts --- but they are the principle behind the pH meter, where a
known reference concentration is compared against an unknown one and the
voltage reports the difference.

\textbf{The check that costs nothing.} $E$ must come out
\emph{positive} if you assigned the anode correctly. A negative answer
means you put the concentrated side as the anode; swap them. And when
the two concentrations become equal, $Q = 1$, $\log Q = 0$ and
$E = 0$ --- the cell dies exactly when the difference disappears.
\end{workedexample}

\begin{yourturn}[YOUR TURN 12 --- four questions]
\begin{enumerate}[leftmargin=*,itemsep=0.5em]
  \item What is $E^{\circ}$ for any concentration cell? Why?
        \workspace[1.6cm]
  \item In a concentration cell, which side is the anode?
        \workspace[1.6cm]
  \item Find $E$ for a copper concentration cell with
        \SI{0.10}{\Molar} and \SI{1.0}{\Molar} ($n = 2$).
        \workspace[2.0cm]
  \item What happens to $E$ as the two concentrations approach each
        other? \workspace[1.6cm]
\end{enumerate}
\selfcheck{(a) 0 --- both half-cells are identical \quad (b) the dilute
side \quad (c) $+0.0296$ V \quad (d) it falls to zero}
\keyonly{\begin{apcanswer}(a) $\mathbf{0}$, because both half-cells
carry the \textbf{same half-reaction}, so their standard potentials are
identical and cancel. (b) The \textbf{dilute} side --- the cell runs so
as to raise the dilute concentration, which requires the electrode there
to dissolve. (c) $Q = 0.10/1.0 = 0.10$, so
$E = -(0.0296)\log(0.10) = -(0.0296)(-1) = \mathbf{+0.0296}$ \textbf{V}.
(d) It \textbf{falls to zero}: when the concentrations are equal
$Q = 1$, $\log Q = 0$, and there is no longer any difference to
exploit.\end{apcanswer}}
\end{yourturn}

% =====================================================================
\section*{\sffamily Ladder 13 \textbullet\ Electrolytic cells}
\zsec{18.7}

Run electricity \emph{into} a cell and you can force an unfavourable
reaction to happen. Same vocabulary, opposite energetics.

\begin{center}
\begin{tikzpicture}[font=\sffamily\footnotesize]
  \node[font=\sffamily\bfseries] at (3.4,3.3) {GALVANIC};
  \node[font=\sffamily\bfseries] at (8.6,3.3) {ELECTROLYTIC};
  \draw[apcgray] (0.5,3.0) -- (11.0,3.0);
  \draw[apcgray] (6.0,3.45) -- (6.0,0.15);
  \foreach \y/\lab/\gal/\ele in {%
      2.5/{$E^{\circ}_{\text{cell}}$}/{positive}/{negative},
      2.0/{$\Delta G^{\circ}$}/{negative}/{positive},
      1.5/{energy}/{\textbf{produces} it}/{\textbf{consumes} it},
      1.0/{driven by}/{the reaction itself}/{an external supply},
      0.5/{oxidation at}/{the \textbf{anode}}/{the \textbf{anode}}} {
    \node[anchor=east,font=\sffamily\bfseries] at (2.0,\y) {\lab};
    \node at (3.9,\y) {\gal};
    \node at (8.6,\y) {\ele};}
  \node[font=\sffamily\footnotesize] at (5.6,-0.1)
    {the last row never changes --- which is why anode/cathode is worth
     learning and $+$/$-$ is not};
\end{tikzpicture}
\end{center}

\begin{workedexample}[I do: forcing a reaction uphill]
\textbf{The zinc--copper reaction has $E^{\circ} = +1.10$ V. What
happens if you apply an external voltage greater than 1.10 V in the
opposite direction?}

\textbf{You drive the reaction backwards.} Copper dissolves and zinc
plates out --- the exact reverse of what the cell does on its own. This
is electrolysis: using electrical energy to force a thermodynamically
unfavourable reaction.

\textbf{Why more than 1.10 V?} You must overcome the cell's own
tendency to run forward. Applying less than 1.10 V lets the cell keep
discharging; applying exactly 1.10 V holds it at balance; applying more
overpowers it.

\textbf{Three examples that matter industrially.} Electroplating
deposits a thin metal layer on an object at the cathode. The
Hall--H\'{e}roult process reduces \ce{Al^3+} to aluminium metal ---
impossible chemically, because aluminium is so easily oxidised, and the
reason aluminium was once more precious than gold. And electrolysing
water gives hydrogen and oxygen, which will not form from water on their
own.

\textbf{What does \emph{not} change.} Oxidation still happens at the
anode and reduction still at the cathode. The electrode \emph{signs}
reverse between cell types, which is precisely why the CED excludes them
and why you should reason with anode/cathode instead.

\textbf{And note the scope boundary.} Zumdahl's \S18.8 covers commercial
electrolysis in detail --- the Downs cell, chlor-alkali, aluminium
production. That section is not on the CED. Understand the
\emph{principle} of electrolysis, which is 9.11; do not memorise
industrial plant designs.
\end{workedexample}

\begin{yourturn}[YOUR TURN 13 --- four questions]
\begin{enumerate}[leftmargin=*,itemsep=0.5em]
  \item What is the sign of $E^{\circ}_{\text{cell}}$ for an
        electrolytic cell? \workspace[1.4cm]
  \item At which electrode does reduction occur in an electrolytic
        cell? \workspace[1.4cm]
  \item Why must the applied voltage exceed the cell's own potential?
        \workspace[1.8cm]
  \item Give one reason aluminium is produced electrolytically rather
        than chemically. \workspace[1.8cm]
\end{enumerate}
\selfcheck{(a) negative \quad (b) the cathode \quad (c) to overpower
the reverse tendency \quad (d) \ce{Al^3+} is too hard to reduce
chemically}
\keyonly{\begin{apcanswer}(a) \textbf{Negative} --- the reaction is not
favoured and must be driven. (b) The \textbf{cathode}, exactly as in a
galvanic cell; only the signs differ, not the definitions. (c) Because
the cell has its own tendency to run \emph{forward}; the applied voltage
must \textbf{overpower} that tendency before the reaction can be pushed
in reverse. (d) Because \ce{Al^3+} has a very negative reduction
potential ($-1.66$ V) --- aluminium is so readily oxidised that no
practical chemical reducing process will reduce it, so electrical
energy must supply the difference.\end{apcanswer}}
\end{yourturn}

% =====================================================================
\section*{\sffamily Ladder 14 \textbullet\ Faraday's law}
\zsec{18.7}

The quantitative side of electrolysis: how much substance a given
current deposits in a given time.

\begin{center}
\begin{tikzpicture}[font=\sffamily\footnotesize]
  \node[font=\sffamily\small\bfseries] at (5.6,2.6)
    {the conversion chain --- never skip a link};
  \foreach \x/\qty in {0.9/{current $\times$ time},
                       3.9/{charge (C)},
                       6.3/{mol \ce{e-}},
                       8.5/{mol substance},
                       10.6/{grams}} {
    \node[align=center] at (\x,1.9) {\qty};}
  \foreach \x/\op in {2.5/{$It$}, 5.2/{$\div F$},
                      7.5/{$\div n$}, 9.6/{$\times M$}} {
    \draw[-{Stealth[length=2mm]},thick] (\x-0.5,1.45) -- (\x+0.5,1.45);
    \node[font=\sffamily\footnotesize,anchor=north] at (\x,1.4) {\op};}
  \draw[apcgray] (0.4,0.75) -- (11.2,0.75);
  \node[anchor=west] at (0.5,0.3)
    {$F = \SI{96485}{\coulomb\per\mole}$ \textbullet\ time must be in
     \textbf{seconds} \textbullet\ $n$ is electrons per formula unit};
\end{tikzpicture}
\end{center}

\begin{workedexample}[I do: electroplating silver]
\textbf{A current of \SI{2.00}{\ampere} runs for \SI{30.0}{\minute}
through a solution of \ce{Ag+}. What mass of silver plates out?}

\textbf{Step 1 --- time to seconds.}
\[ 30.0\;\text{min} \times 60 = \SI{1800}{\second} \]

\textbf{Step 2 --- charge.}
\[ Q = It = (2.00)(1800) = \SI{3600}{\coulomb} \]

\textbf{Step 3 --- moles of electrons.}
\[ \frac{3600}{96485} = \SI{0.0373}{\mole}\;\ce{e-} \]

\textbf{Step 4 --- moles of silver.} The half-reaction is
$\ce{Ag+ + e- -> Ag}$, so one electron per silver:
\[ 0.0373\;\text{mol}\;\ce{Ag} \]

\textbf{Step 5 --- grams.}
\[ 0.0373 \times 107.87 = \mathbf{\SI{4.02}{\gram}} \]

\textbf{The step that separates silver from copper.} For
$\ce{Cu^2+ + 2e- -> Cu}$ you would divide the moles of electrons by
\textbf{two}, because each copper needs two. Getting $n$ from the
half-reaction --- not guessing it --- is the whole skill.

\textbf{Two unit traps.} Time must be in \textbf{seconds}; a current in
amperes is coulombs per second, so minutes or hours must be converted
first. And an ampere-hour is not a coulomb --- multiply by 3600 if a
question gives you one.

\textbf{Sanity check by scale.} Roughly \SI{96500}{\coulomb} deposits
one mole of a singly charged metal. We passed \SI{3600}{\coulomb},
about 1/27 of that, so we should get about 1/27 of a mole of silver ---
roughly \SI{4}{\gram}. It does.
\end{workedexample}

\begin{yourturn}[YOUR TURN 14 --- four questions]
Use $F = \SI{96485}{\coulomb\per\mole}$, $M(\ce{Cu}) =
\SI{63.55}{\gram\per\mole}$.
\begin{enumerate}[leftmargin=*,itemsep=0.5em]
  \item How many coulombs pass in \SI{20.0}{\minute} at
        \SI{1.50}{\ampere}? \workspace[1.8cm]
  \item How many moles of electrons is that? \workspace[1.8cm]
  \item What mass of copper plates out from \ce{Cu^2+}?
        \workspace[2.0cm]
  \item Why must you divide by 2 for copper but not for silver?
        \workspace[1.8cm]
\end{enumerate}
\selfcheck{(a) \SI{1800}{\coulomb} \quad (b) \SI{0.0187}{\mole} \quad
(c) \SI{0.593}{\gram} \quad (d) \ce{Cu^2+} needs two electrons}
\keyonly{\begin{apcanswer}(a) $t = 1200\;\text{s}$, so
$Q = (1.50)(1200) = \mathbf{\SI{1800}{\coulomb}}$. (b)
$1800/96485 = \mathbf{\SI{0.0187}{\mole}}\;\ce{e-}$. (c)
$\ce{Cu^2+ + 2e- -> Cu}$, so $0.0187/2 = \SI{0.00933}{\mole}$ of copper;
$0.00933 \times 63.55 = \mathbf{\SI{0.593}{\gram}}$. (d) Because the
half-reaction $\ce{Cu^2+ + 2e- -> Cu}$ requires \textbf{two electrons
per copper atom}, whereas $\ce{Ag+ + e- -> Ag}$ requires only one. The
divisor comes from the half-reaction, never from
guesswork.\end{apcanswer}}
\end{yourturn}

% =====================================================================
\section*{\sffamily Mixed set \textbullet\ ten questions, no labels}

\begin{apcnote}[Why this section exists]
Every question so far arrived with its ladder attached. These ten do
not --- and in this chapter the first decision is usually \emph{which
relationship}: the table, $E^{\circ}_{\text{cell}}$, $-nFE^{\circ}$,
$(0.0592/n)\log K$, Nernst, or Faraday.
\end{apcnote}

\begin{yourturn}[MIXED SET --- first five]
\begin{enumerate}[leftmargin=*,itemsep=0.5em,label=(\alph*)]
  \item Find the oxidation number of Mn in \ce{MnO2}.
        \workspace[1.6cm]
  \item Using $E^{\circ}$: \ce{Ag+}/\ce{Ag} $= +0.80$,
        \ce{Ni^2+}/\ce{Ni} $= -0.25$ V, find
        $E^{\circ}_{\text{cell}}$ for a \ce{Ni}/\ce{Ag} cell.
        \workspace[2.0cm]
  \item For that cell, $n = 2$. Find $\Delta G^{\circ}$ in kJ.
        \workspace[2.0cm]
  \item Which electrode is written on the left in cell notation, and
        what happens there? \workspace[1.8cm]
  \item A cell has $E^{\circ} = -0.40$ V. Is it galvanic or
        electrolytic? \workspace[1.6cm]
\end{enumerate}
\selfcheck{(a) $+4$ \quad (b) $+1.05$ V \quad (c)
\SI{-203}{\kilo\joule} \quad (d) the anode; oxidation \quad (e)
electrolytic}
\keyonly{\begin{apcanswer}\textbf{(a)} $x - 4 = 0$, so
$\mathbf{+4}$. \textbf{(b)} Silver is higher, so it is the cathode:
$0.80 - (-0.25) = \mathbf{+1.05}$ \textbf{V}. \textbf{(c)}
$-(2)(96485)(1.05) = \SI{-202600}{\joule} \approx
\mathbf{\SI{-203}{\kilo\joule}}$. \textbf{(d)} The \textbf{anode}, where
\textbf{oxidation} occurs. \textbf{(e)} \textbf{Electrolytic} --- a
negative $E^{\circ}$ means the reaction must be
driven.\end{apcanswer}}
\end{yourturn}

\begin{yourturn}[MIXED SET --- last five]
\begin{enumerate}[leftmargin=*,itemsep=0.5em,label=(\alph*)]
  \item A cell has $E^{\circ} = +0.0592$ V with $n = 2$. Find
        $\log K$. \workspace[1.8cm]
  \item Balance $\ce{Cr2O7^2- -> Cr^3+}$ in acid.
        \workspace[2.0cm]
  \item A \ce{Zn}/\ce{Cu} cell ($E^{\circ} = 1.10$ V, $n = 2$) has
        $[\ce{Zn^2+}] = \SI{1.0}{\Molar}$ and $[\ce{Cu^2+}] =
        \SI{0.10}{\Molar}$. Find $E$. \workspace[2.2cm]
  \item \SI{1.00}{\ampere} flows for \SI{96485}{\second} through
        \ce{Ag+}. How many moles of silver plate out?
        \workspace[2.0cm]
  \item In an electrolytic cell, where does reduction occur?
        \workspace[1.4cm]
\end{enumerate}
\selfcheck{(a) 2 \quad (b) $\ce{Cr2O7^2- + 14H+ + 6e- -> 2Cr^3+ +
7H2O}$ \quad (c) $+1.07$ V \quad (d) \SI{1.00}{\mole} \quad (e) the
cathode}
\keyonly{\begin{apcanswer}\textbf{(a)} $\log K = nE^{\circ}/0.0592 =
(2)(0.0592)/0.0592 = \mathbf{2}$, so $K = 100$. \textbf{(b)}
$\ce{Cr2O7^2- + 14H+ + 6e- -> 2Cr^3+ + 7H2O}$; charge checks as
$-2+14-6 = +6$ on both sides. \textbf{(c)} $Q = 1.0/0.10 = 10$, so
$E = 1.10 - (0.0296)(1) = \mathbf{+1.07}$ \textbf{V} --- lower than
standard, because product is in excess. \textbf{(d)}
$Q = (1.00)(96485) = \SI{96485}{\coulomb}$, which is exactly one mole of
electrons; $\ce{Ag+ + e- -> Ag}$ needs one each, so
$\mathbf{\SI{1.00}{\mole}}$ of silver. \textbf{(e)} At the
\textbf{cathode} --- the same as in a galvanic cell.\end{apcanswer}}
\end{yourturn}

% =====================================================================
\section*{\sffamily Mastery tracker}

Tick a row only if \textbf{all four} YOUR TURN questions were right on
the first attempt. Every row is assessed, within the scope stated at the
top: \S18.5--18.6 and \S18.8 are off the CED and are not here, and the
two exclusion statements (electrode $+$/$-$ labelling, agent
terminology) are honoured throughout.

\renewcommand{\arraystretch}{1.25}
\noindent\begin{tabular}{@{}c p{5.7cm} c p{4.9cm}@{}}
\toprule
\textbf{First try?} & \textbf{Skill} & \textbf{Ladder} &
  \textbf{If not, re-read\ldots}\\
\midrule
$\square$ & oxidation numbers & 1 & the sum equals the charge \\
$\square$ & balancing in acid & 2 & O, then H, then charge \\
$\square$ & balancing in base & 3 & add \ce{OH-} to both sides \\
$\square$ & the galvanic cell & 4 & An Ox, Red Cat \\
$\square$ & cell notation & 5 & anode on the left \\
$\square$ & reduction potentials & 6 & do not scale $E^{\circ}$ \\
$\square$ & $E^{\circ}_{\text{cell}}$ & 7 & cathode minus anode \\
$\square$ & spontaneity & 8 & signs run opposite to $\Delta G$ \\
$\square$ & $\Delta G^{\circ} = -nFE^{\circ}$ & 9 & find $n$ properly \\
$\square$ & $E^{\circ}$ and $K$ & 10 & $(0.0592/n)\log K$ \\
$\square$ & the Nernst equation & 11 & solids omitted from $Q$ \\
$\square$ & concentration cells & 12 & $E^{\circ} = 0$, dilute is anode \\
$\square$ & electrolytic cells & 13 & driven, not driving \\
$\square$ & Faraday's law & 14 & seconds, and $n$ per formula unit \\
\midrule
$\square$ & \textbf{mixed set} (8 of 10) & --- & whichever it exposed \\
\bottomrule
\end{tabular}

\begin{apcnote}[Scoring yourself honestly]
14/14 and the mixed set: Unit 9 is complete, and with it the whole
course. Every self-study chapter is now behind you.

10--13: check where the misses fall. Ladders 1--3 are redox bookkeeping
carried over from Chapter 4 and repair with practice. Ladders 6--8 are
sign conventions, and almost every error there is a \textbf{sign} rather
than a misunderstanding --- write ``cathode minus anode'' at the top of
your working every time. Ladders 9--12 are Chapter 17 in new units; if
they missed, the problem may be upstream.

9 or fewer: the likeliest single cause is the opposite sign conventions
of $E^{\circ}$ and $\Delta G^{\circ}$. Re-read Ladders 8 and 9 together
and fix one sentence in your head --- \emph{a battery works on its own
and has a positive voltage, so favoured means $E^{\circ} > 0$ and
$\Delta G^{\circ} < 0$.} Nearly everything else in this chapter follows
from getting that pair the right way round.

\textbf{Where this sits in the course.} Unit 9 combines thermodynamics
(Chapter 17) and electrochemistry (this one) into a single 7--9\% unit.
If both chapters are solid, you have finished the content and what
remains is review.
\end{apcnote}

\end{document}
